% Normal singular comparison kernel; full root and representations written audits passed.
\subsection{A normal singular base with a nonzero generic-comparison kernel}

\begin{lemma}[Proved: residue-field version of the dyadic Bockstein parity]
\label{C:dyadic_residue_Bockstein}
Let $A$ be a discrete valuation ring of characteristic zero whose
maximal ideal is $(2)$, and let $\kappa=A/2A$.
In the exact category of finite free $A/4A$-modules with an
involution $T$, put $D=T-1$ and
\[
 H(M)=\ker(\overline D)/\operatorname{im}(\overline D).
\]
For a cycle $\bar v$, write $Dv=2w$ and define
$N_M[\bar v]=[\bar w]$.
For every monic irreducible $f\in\kappa[t]$ with
$f(t+1)=f(t)$, the function
\[
 \chi_f(M)=\operatorname{length}_{\kappa[t]}H(M)_f\pmod2
                                                               \tag{C.1.N1}
\]
is an additive invariant of the exact representation group.
It vanishes on the trivial and sign lines if $f\ne t,t+1$.
\end{lemma}
\begin{proof}
The proof of Lemma~\ref{C:Bockstein_parity} uses only that
$\operatorname{Ann}_{A/4A}(2)=2A/4A$ and that its quotient is the
characteristic-two field $\kappa$, so it applies here. We spell
out the exact-sequence point. The identity $D^2=-2D$ makes
$\bar D$ a square-zero differential. Division by two determines
$\bar w$ uniquely, changing a lift changes it by a boundary,
and a boundary $\bar v=\bar D\bar z$ maps to a boundary because
$D^2z/2=-Dz$. Thus $N_M$ is well-defined and $\kappa$-linear.

For an admissible extension, choose an underlying free-module
splitting and write
\[
 T_E=\begin{pmatrix}T_L&X\\0&T_Q\end{pmatrix},\qquad
 D_LX+XD_Q+2X=0.
\]
The resulting exact cyclic homology sequence is
$H(L)\to H(E)\to H(Q)\xrightarrow{\partial}H(L)$.
The first two maps commute with $N$, while the displayed
identity gives $N_L\partial=\partial(N_Q+1)$.
The assumption $f(t+1)=f(t)$ therefore preserves the $f$-primary
part, including at the connecting map, and preserves its length.
If the three successive image lengths are $r,s,u$, the three
space lengths are $u+r,r+s,s+u$. Their alternating difference
is even, proving additivity modulo two. On a trivial or sign
line $N$ is respectively zero or one, proving the last claim.
\end{proof}

\begin{theorem}[Proved: normality and trivial generic projective classes do not remove the kernel]
\label{C:normal_trivialK0_kernel}
Let $\zeta$ satisfy $\zeta^2+\zeta+1=0$, put
\[
 A=\mathbb Z_{(2)}[\zeta],\qquad
 R=\left(A[x,y]/(xy-4)\right)_{(2,x,y)},\qquad
 B=R[1/2].                                                \tag{C.1.N2}
\]
Then $R$ is a normal, singular noetherian local domain
of characteristic zero and dimension two, flat over $\mathbb Z_{(2)}$. Every nonzero integer
localization of $R$ has $K_0=\mathbb Z$, with the free module
as generator. In particular, $B$ is a connected regular
principal ideal domain and $K_0(B)=\mathbb Z$.
Nevertheless
\[
 \ker\bigl(R_R(C_2)\longrightarrow R_B(C_2)\bigr)\ne0,
                                                               \tag{C.1.N3}
\]
and $\vartheta_R$ is not surjective in degree two.
\end{theorem}
\begin{proof}
The polynomial defining $A$ is irreducible over $\mathbb Q$
and modulo two. Hence $A$ is a finite local domain over
$\mathbb Z_{(2)}$, its unique maximal ideal is $(2)$,
and its residue field is $\mathbb F_4$.
It is a discrete valuation ring by the principal-maximal-ideal
criterion \cite[Lemma 10.119.7, Tag 00PD]{C-NormalKernelStacks}.
Write $\alpha$ for the image of $\zeta$ in $\mathbb F_4$.

The polynomial $xy-4$ is primitive and irreducible as a
linear polynomial in $y$ over the UFD $A[x]$: its coefficients
$x$ and $-4$ have greatest common divisor one.
Thus $R$ is a noetherian local domain of characteristic zero.
It is torsion free, hence flat, over the DVR $A$ (and over
$\mathbb Z_{(2)}$). It is a hypersurface in the three-dimensional
regular local ring $A[x,y]_{(2,x,y)}$, so it has dimension two
and is Cohen--Macaulay. In particular it satisfies $(S_2)$.
The relation $xy-4$ belongs to the square of the ambient
maximal ideal. Thus $2,x,y$ remain linearly independent in
$\mathfrak m_R/\mathfrak m_R^2$, and $R$ is singular.

After inverting two, the equation identifies $y$ with $4/x$.
Consequently $B$ is a localization of
$\operatorname{Frac}(A)[x,x^{-1}]$, a principal ideal domain.
It is a nonzero domain and is regular. A height-one prime of
$R$ containing two is either $(2,x)$ or $(2,y)$, because
$R/2R=\mathbb F_4[x,y]_{(x,y)}/(xy)$.
At $(2,x)$ the element $y$ is a unit, so $x=4/y$ and the
maximal ideal is generated by two. This localization is a DVR;
the other one is identical with $x$ and $y$ interchanged.
Every other height-one localization is regular because it
is a localization of $B$. Hence $R$ satisfies $(R_1)$, and
Serre's criterion gives normality
\cite[Lemma 10.157.4, Tag 031S]{C-NormalKernelStacks}.

Finite projectives over a local ring are free, and finite
projectives over a principal ideal domain are free. Thus the
rank identifies both $K_0(R)$ and $K_0(B)$ with $\mathbb Z$.
An odd integer was already a unit of $R$, while inverting an
even nonzero integer is the same as inverting two. This proves
all the stated integer-localized $K_0$ assertions.

Put $a=1+2\zeta$. Since $a^2=-3$, the matrix
\[
 T=\begin{pmatrix}a&x\\y&-a\end{pmatrix}               \tag{C.1.N4}
\]
satisfies $T^2=(a^2+xy)I_2=I_2$ over $R$.
Let $M=(R^2,T)$ and set
$\eta=[M]-[\mathbf1]-[\mathrm{sgn}]$.
Over $B$ the matrix
\[
 P=\begin{pmatrix}x&x\\1-a&-1-a\end{pmatrix}
\]
has determinant $-2x$, a unit, and satisfies
$TP=P\operatorname{diag}(1,-1)$. Therefore $\eta$ maps to
zero in $R_B(C_2)$; its vanishing does not merely rely on a
rank computation.

There is a local ring homomorphism
\[
 R\longrightarrow A/4A,\qquad x\longmapsto2,\quad
                                  y\longmapsto2\zeta.     \tag{C.1.N5}
\]
The equation is respected because $4\zeta=4=0$ in the target.
Every inverted denominator has nonzero residue in
$\mathbb F_4$, so the map is defined on the stated localization.
For this specialized involution $\bar D=0$ and
\[
 N_M=\begin{pmatrix}\alpha&1\\\alpha&1+\alpha\end{pmatrix},
 \qquad\det(tI-N_M)=t^2+t+\alpha^2.                      \tag{C.1.N6}
\]
The image of $u\mapsto u^2+u$ on $\mathbb F_4$ is
$\{0,1\}$. Since $\alpha^2\notin\{0,1\}$,
$f=t^2+t+\alpha^2$ is irreducible. It also satisfies
$f(t+1)=f(t)$. Thus $H(M)$ is a simple
$\mathbb F_4[t]/(f)$-module and
$\chi_f(M)=1$. The trivial and sign lines have invariant zero.
Base change is exact on the representation category since
all admissible sequences split on their underlying projectives.
Lemma~\ref{C:dyadic_residue_Bockstein} therefore proves
$\eta\ne0$.

Finally the degree-two polynomial image is generated by the
trivial line and the regular permutation module. The admissible
sequence
$0\to\mathbf1\to R[C_2]\to\mathrm{sgn}\to0$
identifies this image with the span of the two lines, both
annihilated by $\chi_f$. The displayed representation is
outside that image. This proves all claims, without asserting
an additive order for $\eta$.
\end{proof}

This example separates the generic-comparison kernel from the
nontrivial-generic-Picard obstruction in
Proposition~\ref{C:normal_domain_counterexample} even among
normal local domains. Its special arithmetic fibre is a singular
node. The argument does not assert a generic $K_0$ computation
for the distinct complete ring
$\mathbb Z_2[[x,y]]/(xy-4)$.

The exact bounded certificate
\texttt{code/C\_finish\_normal\_kernel.py} verifies the defining-ring
matrix identity, the explicit generic eigenbasis, the four-element
field irreducibility test and the specialized primary invariant.
It also checks all $2304$ candidate upper blocks in the line/line
and witness/line extension ranges: all $240$ admissible extensions
satisfy the predicted parity. All $180$ changes of basis in
$\mathrm{GL}_2(\mathbb F_4)$ preserve the primary calculation.
Its scope and runtime are recorded in
\texttt{logs/C\_finish\_normal\_kernel.json}; the arbitrary-rank
additivity and commutative-algebra assertions have the independent
written proofs above.
